\honest{Something to Protect}{Eliezer Yudkowsky}{2008}

In Japanese fiction, power comes from having someone to protect. In the \textsc{X} saga each hero draws power from one most precious person, and if that person is killed or hurt the wrong way, the hero loses it, from despair, as good as being taken off the board. In Western comics, the hero ``gets bitten by a radioactive spider'' and then needs something to do with his powers, ``so he decides to fight crime,'' and then whines about how much time it takes. In Western comics power comes first and purpose after; in Japanese fiction, often, the reverse. \nb{That hero is Spider-Man. He fights crime because he let a burglar escape and the burglar killed his uncle. His motive is guilt, not a need for something to do.}

Unhappy Westerners are told to pick an altruistic cause that suits them, ``like picking out nice living-room drapes,'' though not too expensive.

``Of course I'm not saying all this to generalize from fictional evidence.'' It is there so that you will not confuse my idea with the Western version.

My idea is that a rationalist must value something more than rationality, or the Art collapses into infinite recursion; and I ask where rationalists come from and how they gain their powers. As the \textsc{Twelve Virtues} says, if your conception of rationality is to believe the Great Teacher, who says the sky is green, saying it is your duty to be rational only enshrines the mistake. Curiosity is as old as humanity, but campfire tales satisfy it just as well.

Science beat authority ``because it displayed greater raw strength in the form of technology, not because science sounded more reasonable.'' Magic and scripture still sound more reasonable to untrained ears, which is why the tension continues. I cite no historian. \nb{One who counted, B. Zorina Khan, found that Britain's great inventors were mostly not scientists until late in the nineteenth century. That does not refute the claim, but the claim needs an argument.} To people who say truth should matter more than usefulness: ``Forget for a moment'' that people ``in pretty much that frame of mind defended the Bible.'' \nb{The post asks the reader to forget the comparison, and does not take it back.} Propositional morality has too many degrees of freedom.

Then I explain that my love of truth is ``more complicated.'' No one masters the art without caring about truth, as few master composers hate music. But I like the discipline of making beliefs yield predictions, and loving true-seeming ideas while ready to drop them. I also like ``declaring that I value mere usefulness above aesthetics,'' which is ``almost a contradiction, but not quite.'' But never deliberately believe a useful falsehood. Morality and aesthetics alone, the belief that one ought to be rational, would not have got humanity out of the authority-hole.

My main argument is a dilemma. Option 1 saves 400 lives for certain. Option 2 saves 500 with probability 90\% and none otherwise. If your daughter is one of the 500, option 1 saves her with probability 80\% and option 2 with 90\%, and love may push you to ``shut up and multiply.'' Even one of the 500 might grandstand and cling to certainty, since our own lives often matter less to us than a good intuition. And everyone in the crowd is someone's child, so altruists too should pick the second. The point is not that one life outweighs 499, but that more than your own life must be at stake before people turn to math. \nb{Framed this way, neither option is certain for her, so the ``comforting feeling of certainty'' the post says love must overcome is gone before love does anything. Her chances, 80 and 90, are the expected numbers saved, 400 and 450, divided by 500. So caring about her and counting everyone always give the same answer, as the post itself says. The example cannot show that love leads anywhere that counting does not. Nor does the post give evidence that love makes people better with probabilities. Experiments by Rottenstreich and Hsee suggest that strong feeling makes people notice a difference like 80 against 90 less, not more.}

Many think rationality means choosing only methods certain to work; but hopefully you care more about your daughter than about ``rationality.'' Pride in being a rationalist can stop you from learning: ``You will only be able to learn something about rationality if your daughter's life matters more to you than your pride as a rationalist.''

Standing alone against a crowd is ``much scarier than a mere threat to your life, according to the revealed preferences of teenagers who drink at parties and then drive home.''

I then tell my own history. I grew up an ``above-average Traditional Rationalist a la Feynman and Heinlein, and nothing more,'' and grew further only once I had ``something terribly important'' to do. I do not say what it was. A few paragraphs later I state: ``No one masters the Way until more than their life is at stake.'' \nb{The only evidence given is this one case.}

You may learn from such a moment if you can say ``I must have had the wrong conception of rationality,'' not ``Look at how rationality gave me the wrong answer!''; mastering rationality takes rationality to bootstrap. A self-image as someone who faces harsh truths helps, but may make it too hard to admit you have done rationality all wrong, as with one who must admit the Great Teacher was a fraud.

Only when you care more about success than any technique do you understand Musashi: ``the Way of the Ichi school is the spirit of winning, whatever the weapon and whatever its size.'' You cannot just pick a cause because you need a hobby. Look for a good cause and your mind supplies a cliché. But ``Learn how to multiply, and perhaps you will recognize a drastically important cause when you see one.''

Early in the post, the Art of rationality ``must have a purpose other than itself.'' At the end, to subordinate its aesthetics to a higher cause ``is part of the aesthetic of rationality,'' and you should appreciate its beauty ``for its own sake.''
